% Created 2023-04-17 Mon 09:02
% Intended LaTeX compiler: pdflatex
\documentclass[pressentation,10pt,aspectratio=169,xcolor=table, colorlinks=true]{beamer}
\usepackage[utf8]{inputenc}
\usepackage[T1]{fontenc}
\usepackage{graphicx}
\usepackage{longtable}
\usepackage{wrapfig}
\usepackage{rotating}
\usepackage[normalem]{ulem}
\usepackage{amsmath}
\usepackage{amssymb}
\usepackage{capt-of}
\usepackage{overlock}\usepackage[english]{babel}
\usepackage{tikz, pgfplots}\usepackage{smartdiagram}
\usetikzlibrary{arrows,shapes,positioning,mindmap,decorations.pathreplacing,backgrounds,overlay-beamer-styles,calc,3d,fit}
\pgfplotsset{/pgf/number format/assume math mode=true}\usepackage{csquotes}
\usepackage[citestyle=verbose, backend=biber]{biblatex}\addbibresource{refs.bib}
\usepackage[type={CC}, modifier={by-sa}, version={4.0},]{doclicense}
\pgfplotsset{compat=1.17}
\usepackage{minted}\usemintedstyle{emacs}\usepackage[]{tcolorbox}
\newtcolorbox[%
auto counter,number within=section]{work}[1][]{colback=black!5!white,colframe=black!50!white,fonttitle=\sffamily\bfseries,title=Work~\thetcbcounter: #1}

\hypersetup{
  colorlinks = true
}
\AtBeginSection[]{%
  \begin{frame}<beamer>
    \frametitle{Outline}
    \tableofcontents[sectionstyle=show/shaded,subsectionstyle=hide]
  \end{frame}
  \addtocounter{framenumber}{-1}% If you don't want them to affect the slide number
}
% Short beamer templating
%% \titlegraphic{\includegraphics[width=0.125\linewidth]{logo_cesbio_transp.pdf}}
\usetheme{default}
\usecolortheme[snowy]{owl}
\setbeamertemplate{navigation symbols}{}{}
\setbeamertemplate{footline}[frame number]
\setbeamertemplate{itemize item}{$\bullet$}
\setbeamertemplate{itemize subitem}{$\star$}
\setbeamercolor{block body}{parent=normal text,use=block title,bg=black!5,}
\setbeamercolor{block title}{fg=black, bg=gray}
\setbeamertemplate{blocks}[rounded][shadow=false,] 

\author[M. Fauvel]{Mathieu FAUVEL}
\date{\today}
\title{UE Apprentissage SIA S9}
\subtitle{Supervise Learning}
\institute{CESBIO, Universit{\'e} de Toulouse, CNES/CNRS/INRAe/IRD/UPS, Toulouse, FRANCE}
\begin{document}
\maketitle

\begin{frame}
  \frametitle{Outline}
  \tableofcontents[hideallsubsections]
\end{frame}

\section{Introduction}

\begin{frame}
  \frametitle{Myself}
  \begin{itemize}
  \item Contact: mathieu.fauvel@inrae.fr
    
  \item CV:
    \begin{itemize}
    \item [2004-2007] Ph.D. degree in Signal and Image Processing from the INP, Grenoble \& the University of Iceland
    \item [2007-2008] Assistant Professor Grenoble
    \item [2008-2010] Post-doc position at INRIA - MISTIS Team
    \item [2010-2011] Assistant Professor Toulouse
    \item [2011-2018] Associate Professor at DYNAFOR \& INP, Toulouse
    \item [Since 2018] Research (CR1) at CESBIO, INRAe
    \end{itemize}
  \item Research interests are:
    \begin{center}
      Machine learning for environmental/ecological monitoring
    \end{center}
  \end{itemize}
\end{frame}

\begin{frame}
  \frametitle{Course Objectives}
  TBC
\end{frame}

\begin{frame}
  \frametitle{What is Supervised (Machine) Learning ?}
  \begin{columns}[t]
    \begin{column}{0.5\linewidth}
      \begin{itemize}
      \item <1-> \textbf{Artificial Intelligence}
        \centerline{\em Perform human tasks using computers and algorithms}
      \item <2-> \cite{mitchelllearning}:
        \begin{center}
          \emph{Machine Learning is defined as the capacity of a computer program to improve its performance measure with observations}
        \end{center}
      \item <3-> \textbf{Deep Learning}
        \begin{itemize}
        \item Lot of data data
        \item Complex model
        \item High performance computing
        \end{itemize}
      \end{itemize}
    \end{column}

    \begin{column}{0.5\linewidth}
      \begin{center}
        \begin{tikzpicture}[scale=0.5]
          \draw[fill=yellow!40] (2.75,0) ellipse (6.4 and 3.2);
          \draw[fill=blue!40] (1,0) ellipse (4 and 2);
          \draw[fill=red!30] (-1,0) ellipse (1.6 and 0.8);
          
          \draw (4,2.3) node {\textbf{\small Artificial Intelligence}};
          \draw (1.7,1.1) node {\textbf{\small Machine Learning}};
          \draw (-0.95,0) node {\small \textbf{DL}};
        \end{tikzpicture}
      \end{center}
    \end{column}
  \end{columns}
\end{frame}


\begin{frame}
  \frametitle{Main notations}
  \begin{block}{Data}
    \begin{itemize}
    \item Observed data \(\mathbf{x}\in \mathbb{R}^d\), called \emph{input variables} or \emph{predictors}.
    \item Data to be predicted \(y\in \mathbb{R}^p\), called \emph{output variables} or \emph{responses}.
    \end{itemize}
  \end{block}

  \begin{block}<2->{Learning Problem}
    \begin{itemize}
    \item If \(y\) are quantitative data: \emph{regression} and \(y\in\mathbb{R}^p\).
    \item If \(y\) are categorical data:  \emph{classification} and \(y\in\{C_1, \ldots, C_C\}\) with \(C_i\) refers to the \(i^{th}\) category.
    \end{itemize}
  \end{block}

  \begin{block}<3->{Prediction function}
    \begin{eqnarray*}
      f_{\boldsymbol{\theta}}: \mathbb{R}^d & \to & \mathbb{R}^p\\
      \mathbf{x}&  \mapsto & y
    \end{eqnarray*}
  \end{block}
\end{frame}

\begin{frame}
  \frametitle{Online References}

  \begin{itemize}
  \item \cite{zhang2021dive}
    
  \item \cite{prince2023understanding}

  \item \cite{pml1Book}

  \item \cite{pml2Book}
  \end{itemize}
  
\end{frame}

\section{Introductory Example}
% Predict function
% Loss and Objective functions
% Gradient
% Optimization and Batch optim

\subsection{Linear model for regression}

\begin{frame}[fragile]
  \frametitle{1D linear model}
  \begin{columns}
    \begin{column}{0.5\linewidth}
      \begin{itemize}
      \item  \(f(x) = wx +b\)  and \(\boldsymbol{\theta} = (w,b)\)     
      \item Loss function: \(\ell \big(f(x_i),y_i\big) = \big(f(x_i) - y_i\big)^2\)
      \item Objective function: \[\mathcal{L} = \frac{1}{n} \sum_{i=1}^n \ell_i\]
      \item Gradients update: \(\boldsymbol{\theta}^{t+1} = \boldsymbol{\theta}^{t} - \eta\nabla_{\boldsymbol{\theta}}\mathcal{L}\)
        \begin{eqnarray*}
          \frac{\partial \mathcal{L}}{\partial w} & = & \frac{2}{n}\sum_{i=1}^n x_i(wx_i +b-y_i)\label{eq:lin:grad:w}\\
          \frac{\partial \mathcal{L}}{\partial b} & = & \frac{2}{n}\sum_{i=1}^n (wx_i +b-y_i)\label{eq:lin:grad:b}
        \end{eqnarray*}
      \end{itemize}
    \end{column}
    \begin{column}<2>{0.5\linewidth}
      \begin{work}
        Implement the 1D regression in pytorch of the following function:
        \[f(x)=\frac{1}{2}\exp(0.1x) + \frac{1}{2}\sin\big(\frac{2\pi}{10}x\big)\]

        Do the notebook: \textcolor{ornage}{regression_1d_toy.ipynb}
      \end{work}
    \end{column}
  \end{columns}
\end{frame}

\subsection{Linear model for classification}

% \section{}

\section{Multi layer Perceptron}

\section{Convolutional Neural Networks}

\section{Transformer}



\begin{frame}[fragile,allowframebreaks,label=]{Bibliography}
\printbibliography
\end{frame}

\begin{frame}[label={sec:orgc74e2e4},noframenumbering]{}
\begin{center}
\small
\doclicenseLongText

\doclicenseImage
\end{center}
\end{frame}

\end{document}
%%% Local Variables:
%%% mode: latex
%%% TeX-master: t
%%% End:
